% ID: SEL-D0040-OB-E02
% ESTADO: cotejada. Segunda pasada visual y borrador compilado revisados.
% Fuente: Selectividad/matematicas_aplicadas/Ua5MATEMÁTICAS APLICADAS examen.pdf
% SHA-256: ac80a164e049669cd01083ca5adaddfa90efbf556b050de2d99f294ab55670c9
% Página: 2; opción B; ejercicio original 2.

Se considera la función $f(x)=x^3-2x^2+x$.
\begin{enumerate}[label=\alph*)]
\item \textbf{(1.3 puntos)} Halle el máximo, el mínimo y el punto de inflexión de la función.
\item \textbf{(0.6 puntos)} Calcule los puntos de corte con los ejes.
\item \textbf{(0.6 puntos)} Obtenga las ecuaciones de las rectas tangentes a la gráfica de $f$ en los puntos de abscisas $x=0$ y $x=1$.
\end{enumerate}
